\chapter{自由法则与自然的统治}

在《法哲学原理》的序言中，黑格尔把这部著作中阐述的“伦理世界”理论与自然哲学进行了对比。在涉及自然的时候，大家都承认，哲学要如其所是地认识自然，自然“自在地”是理性的，并且科学要概念把握自然中的现在的、现实的理性，然而由人所产生的且依赖于人的意志的伦理—历史世界则处于任意和偶然的摆布之下，就此而言，就应该是无法则的和无理性的。\footnote{《法哲学原理》，序言，《全集》第7卷，H. 格洛克纳编，1927-1930年，第23页以下（《百年纪念版》）。}

关于自然规律和法权法则之间的区别这一哲学史古老主题，黑格尔自己并没有在《法哲学原理》的文本本身中，而是在柏林时期的讲座导言中进行了处理。而在甘斯出版的第二版《法哲学原理》（1833年）的序言中，紧接在这个对比之后，则有一个较长的补充。\footnote{甘斯编辑的注释以H. G. 霍托1822-1823年冬季学期讲座笔记中的一段话为基础，参见《法哲学：依据黑格尔教授先生演讲（1822-1823年）》（Philosophie des Rechts, Nach dem Vortrage des H. Prof. Hegel, im Winter 1822/1823），柏林，手写笔记第2册，第1-4页（柏林普鲁士文化遗产基金会国家图书馆黑格尔遗稿）。进一步参见《1824-1825年冬季学期法哲学讲演录》（Vorlesungen über Philosophie des Rechts WS 1824/1825），格里斯海姆笔记，四开本德语手稿第545号，第13页以下。[这两本笔记现收入《黑格尔全集》（历史考订版）第26卷第2分册和第3分册（G. W. F. Hegel, Gesammelte Werke, Band 26, 2-3. hrsg. von der Rheinisch-Westfälischen Akademie der Wissenschaften, Felix Meiner Verlag Hamburg 2015）。——译者]}按照这里所做的区分（这种区分同样依据一种古老的概念分类），有“两类规律或法则(Gesetzen)，自然规律和法权法则”。自然规律确实存在，它们的存在就使它们有效，由此我们即把“一般自然”思维为由规律所规定的。存在(Sein)和效力（应当）在自然规律领域是一致的：“因为这些规律是准确的，只有我们对这些规律的观念才会错误。这些规律的尺度是在我们身外的，我们的认识对它们无所增益，也无助长作用，我们对它们的认识可以扩大我们的知识领域，如此而已。”\footnote{参见《法哲学：依据黑格尔教授先生演讲（1822-1823年）》，第1页，以及《法哲学原理》序言，第24页，附释。（中译文引自《法哲学原理》，第14页。——译者）}就此而言，自然规律与法权法则没有区别，我们也是从“外部”、作为某种给定的东西认识法权法则的，而法律人也必然停留在这种给定存在的立场上。按照这种立场，法则(Gesetze)之所以有效，乃是因为它们是被设定起来的(gesetzt)。然而，从法权法则的差异性和可变性就可以非常明显地看出，它们并不像自然规律那般“绝对”，因此存在和效力在这里就是彼此分离的。法权法则在字面意思上是“被设定的东西”，一种“源出于”人的东西，并依赖于人的意志和意识。自然规律直接规定物的存在，其效力根本而言是不可变的，法权法则由于这种依赖性而处于不断的变化之中。法权法则具有历史，因此人不会让它们直接存在和有效；相反，它们的存在和效力都要以人自身即人的意志和意识的主体性为中介。我们无法像服从自然规律的必然性那样服从由一位立法者所颁布的法律的权力和权威的必然性，因为自然规律调整的乃是人之外的物的运行，而法权法则的尺度不再在我们之外，而是就在我们之中。

在自然界中，一条规律的最高确证是，它一般地存在。而在法权领域，一条法则的确证并不是在其存在中，而在于它被认知和意求，因为只有在这里存在和应当之间的冲突才是可能的。在自然规律和法权法则的那种对比中，最重要的立场并不是自霍布斯、休谟和康德以来出现在自然法基本概念中的二元论。黑格尔感兴趣的是这种概念体系中更深层次上的二元论——自然与精神之间的对立，它通过激化存在和应当的对立来扬弃这一对立。黑格尔的对比进一步化为这种尖锐的对立：

\begin{quote}
如果我们对这两类法则（或规律）的这种差别进行考察并追问，法权法则所归属的基地是什么，那么我们就会看到：法只产生于精神，因为自然没有任何法。因此就有一个存在着的自然世界和一个与自然相对立的精神世界。\footnote{《法哲学：依据黑格尔教授先生演讲（1822-1823年）》，第2页。（本段引文见《黑格尔全集》（历史考订版）第26卷第2分册，第772页。——译者）}
\end{quote}

在这种意义上，黑格尔在《法哲学原理》第4节中说，法的基地一般说来是“精神的东西”，并且其出发点是自身意求着的“自由意志”，自由意志在法的体系中有其历史形态、“实现了的自由的王国”，即“从精神自身产生出来的、作为第二自然的那个精神的世界”\footnote{《法哲学原理》，导论，第4节，第50页。（中译文参见《法哲学原理》，第10页。——译者）}。

作为第二自然的精神世界，并不是亚里士多德的那种以礼法($\nu\acute{o}\mu o\varsigma$)和习惯($\check{\epsilon}\theta o\varsigma$)为基础的素朴伦常和城邦伦理的第二自然($\delta\epsilon\acute{\upsilon}\tau\epsilon\rho\alpha\ \phi\acute{\upsilon}\sigma\iota\varsigma$)，而是一种由人产生并设定于作品之中的自然。就此而言，较之于亚里士多德的《尼各马可伦理学》或《政治学》，这种第二自然就更接近于霍布斯的《利维坦》。霍布斯在现代自然法理论的开端处即借助“自然”社会和“人为”社会的经典对比，削弱了亚里士多德的城邦的自然理论以及传统的智者派的礼法和自然的二分理论。实际上，几乎所有黑格尔借以区分自然规律和法权法则的那些谓词，都可以在这种经典对比中重新找到。“自然社会”中的和谐（霍布斯和黑格尔都认为，这种和谐是动物式的和谐）是“以自然为中介的神的作品”；与之相反，人的和谐则是“人为的”，它以意志的协调一致（契约）为中介。\footnote{霍布斯：《法律原理》（Elements of Law）第1卷，第19章，第5节；《论公民》（De cive）第2卷，第5章，第5节；《利维坦》第2卷，第17章。}与黑格尔一致，霍布斯首先也把自然物视为是持存的和必然的东西，而人所产生的法律和契约则是人为的作品，由于依赖于人的意志，它们便被视作任意的和可变的。个中缘由完全在于自然(Physis)和礼法(Nomos)这一对概念主题：自然性的东西不依赖于契约和行为（在黑格尔看来，则是不依赖于意志和意识），而人为的东西则始终受制于它们。但是霍布斯从这一对比中得出了一个不同于智者的二分以及柏拉图和亚里士多德对这种二分消解的结论。霍布斯不再想要表明，唯有一个自然社会（也就是以自然尺度为准的社会）才能对个体施加义务，一个人为的社会则让个人保持自由；相反，他想要证明，唯有一个严格意义上的“人为”社会才可以向公民施加服从的义务——在涉及自然的时候，已经不再能够有意义地使用“义务”概念了。就此而言，与礼法和自然的二分相对，义务的思想恰好从“自然”颠倒为“人为”。\footnote{对此，参见沃尔夫（F. O. Wolf）：《霍布斯的新科学：论近代政治哲学的基础》（Die neue Wissenschaft des Thomas Hobbes. Zu den Grundlagen der politischen Philosophie der Neuzeit），斯图加特-巴德·坎施塔特，1969年，第82页以下。}

这种颠倒对近代自然法理论的意义，可以从圣经中造物主的功能上看出来：造物主把自然性的东西贬为一种被造物，自然性的东西不再在自身中，而是在上帝的意志有其根据。对于义务的概念来说，这就意味着：如果现在还要宣布自然性的东西对人具有约束力，那么这种约束力的根据就不再是自然，而只在与自然相分离的上帝的意志之中。这种模式引起了更进一步的转变：一种完全不同的约束力与一般“自然”的约束力相对立，同时也抛弃了从圣经的造物主的意志中推导出这种约束力的做法。\footnote{参见上书，第84页。沃尔夫明确地阐述了这一观点。}在黑格尔看来，迈出这一步的不是霍布斯而是卢梭。他在《法哲学原理》第258节中把这算作卢梭的功劳：“一个原则，它不仅在形式上（诸如社会性冲动、神的权威），而且在内容上也是思想，并且是思维本身。这就是说，他提出意志作为国家的原则。”\footnote{《法哲学原理》，第258节，附释，第330页。（中译文参见《法哲学原理》，第254页。——译者）}卢梭把“人的最内在的东西”、作为“与自身相统一的”自由提升为法的基础，并且借此给了人以对抗自然或者说转移到自然上的上帝意志的“一种无限的力量”。\footnote{《哲学史讲演录》第3卷，《全集》第19卷，第527页。（中译文参见〔德〕黑格尔：《哲学史讲演录》第四卷，贺麟、王太庆译，商务印书馆1978年版，第233页。——译者）}同样，自然、自由与法之间的对立已经出现在霍布斯的作品中。在黑格尔看来，现代自然法的真正历史开始于霍布斯，因为霍布斯最先“试图把维系国家统一的力量、国家权力的本性回溯到内在于我们自身的原则，亦即我们承认是我们自己所有的原则”。\footnote{参见《哲学史讲演录》第3卷，《全集》第19卷，第442页。（中译文引自《哲学史讲演录》第四卷，第157页。——译者）}霍布斯与这种假设决裂：存在着一种超出人的意志之上的、支配着万物的自然目的法则(lex naturalis)，理性要成为“正确的理性”(recta ratio)并使“正义的”行动成为可能，就必须摹写、效仿这种法则。在霍布斯那里，就自然的“规律”包含了使得抛弃自然法成为必然条件而言，法和社会的理论仍居于自然的条件之下。通过霍布斯，在自然法的传统名称中产生了一种“歧义”，对黑格尔来说，这种歧义至关重要，他在霍布斯讲座中也把注意力放在了这种歧义上面：“‘自然’这一名词具有歧义，一方面，人的自然指他的精神性、合理性而言；另一方面，他的自然状态则指人按照他的自然性而行动的状态而言。”\footnote{同上书，第443页。（中译文引自《哲学史讲演录》第四卷，第159页。译文稍有改动。——译者）}在黑格尔看来，卢梭加剧了霍布斯所进行的与传统自然法理论的决裂，这种歧义在卢梭那里进一步升级。卢梭把人的“精神性、合理性”理解为人的自由，自由把人从自然中提升出来，并且成为区分人与动物的根本标志。霍布斯在基于契约的制度性的或者“人为的”统治团体之外，还让主奴关系的“自然性的”统治关系继续存在下去，而卢梭则一般地追问统治的绝对合法性。卢梭既不是将实定法，也不是将霍布斯作为基础的自然规律，而是将自由作为合法性的原则。卢梭并不考虑他的17世纪先行者，在他看来，他们与为奴役人、压迫人进行辩护的亚里士多德和实施这种行径的卡利古拉同属一类。同时，卢梭也不考虑欧洲国家现行有效的法权法则：“他对上述问题做出这样的回答：人是有自由意志的，因为‘自由是人的品德。放弃自己的自由，就是放弃做人。放弃自由，就是放弃一切义务和权利’。”\footnote{《哲学史讲演录》第3卷，《全集》第19卷，第527页。（中译文引自《哲学史讲演录》第四卷，第233页。——译者）黑格尔引用了《社会契约论》第一卷第四章的著名段落，参见第一章、第七章、第十一章；第三卷，第九章，注释；《论人类不平等的起源》（Discours sur l'origine de l'inégalité parmi les hommes），K. 维甘德编，汉堡，1955年，第106页以下。另参见《世界历史哲学讲演录》第4卷，第920页以下。}诚如黑格尔所言，卢梭的原则：人是自由的，并且建立在“公意”基础上的国家是这种自由的实现，“乃是正确的”。在卢梭这里，只有当他让公意由众多个体的意志、它们的自然性的自由倾向来“构成”的时候，才出现了“歧义”。由此，“自然的自由”——个体意志的自发性——又取消了“作为完全绝对物的自由”，卢梭在其公意概念中思考的正是这种自由。\footnote{参见上书，第528页：“这些原则表达得很抽象，我们必须把它们正确地找出来；可是歧义立即就发生了。人是自由的，这当然是人的实质本性，这种本性在国家里不但没有被扬弃，事实上倒是开始被建立起来了。本性的自由，自由的禀赋并不是现实的；因为国家才是自由的实现。”（中译文引自《哲学史讲演录》第四卷，第234页。——译者）此外，参见《哲学全书》，第163节，《全集》第8卷，第360页；《法哲学原理》，第258节，第530页以下。}同时，卢梭也使我们意识到，自由乃是完全的绝对物、“人的概念”：

\begin{quote}
恰恰就是思维本身；要是抛开思维来谈自由，就不知道自己说的是什么东西。思维的自身统一性就是自由，就是自由意志……意志只有作为思维的意志才是自由的。现在自由的原则出现了，它把这种无限的力量给予了把自己理解为无限者的人。\footnote{参见《哲学史讲演录》第4卷，《全集》第19卷，第528页以下。（中译文引自《哲学史讲演录》第四卷，第234页。——译者）}
\end{quote}

意志于其自身之所是，既不能通过与自然及其“规律”的类比，也不能通过意志与自然或自然性地得到规定的自由任意之间的单纯对立来把握。意志——卢梭诉诸“公意”，作为国家的基本原则，使我们想到这一点——必须超出这些对立。意志只有在这种情况下才是自由的，即它“意求的不是他者、外在的东西和陌生的东西，因为如果这样它就是依赖性的，而只是意求自身——意求意志。绝对意志就是意求成为自由”。卢梭用这种自由意志的理念，完成了由霍布斯开启的意求和思维自身的个体的主体性对[传统]自然目的论的颠覆。这种自然目的论构成了传统自然法的自然法则(lex naturalis)的基础。取代自然的约束力以及霍布斯对造物主意志的如此模糊的诉求的，乃是绝对化了的意志的约束力。黑格尔认识到了卢梭对霍布斯模式的这种转变，并且承认这一转变乃是他自己的法概念的最重要的前提：

\begin{quote}
绝对意志就是意求成为自由。意求自身的意志乃是一切权利和一切义务的根据，因此也就是一切法权法则、义务要求和所施加的约束力的根据。意志自由本身就是一切法的原则和实体性基础，它自身就是绝对的、自在自为的永恒的法和最高的法。就此而言，其他的、特殊的法被置于次要的地位，它甚至是使人之成为人的东西，因此也就是精神的基本原则。\footnote{《世界历史哲学讲演录》第4卷，第921页以下。}
\end{quote}

由此出发，黑格尔推出了向康德哲学的过渡。康德通过将自然的立法与自由的立法、经验的意志与“自由的和纯粹的意志”彻底分离开来，为终结在卢梭那里尚且存在的自然法概念中的“歧义”做好了准备。\footnote{《哲学史讲演录》第4卷，《全集》第19卷，第529页以下，第552、590页。}在先验哲学的立场上，意志的“概念”达到了对其自身的理解。“自我意识的简单统一”、“自我”是从一切给定的自然秩序中摆脱出来的、“密不透风的、完全独立的自由和一切普遍规定（也就是思维规定）的渊源，——理论理性，同样也是一切实践规定的最高的东西，——作为自由的和纯粹的意志的实践理性；而意志的理性在于，自身保持在纯粹的自由之中，在一切特殊物中只意求这种纯粹自由，仅仅为了法而意求法，仅仅为了义务而意求义务”。\footnote{《世界历史哲学讲演录》第4卷，第922页。}

由此，“对精神性的东西的意识”将成为法哲学的基础；它发现了一种“国家的思想原则，现在，这种原则不再是诸如社会性冲动、财产安全的需要之类的任何一种意见的原则，也不是像君权神授这样的虔诚原则，而是与我的自我意识同一的确定性原则[……]”\footnote{同上书，第924页。参见拙文，《黑格尔对自然法的批判》，《黑格尔研究》第4卷，1967年，第178页以下。现在也收入《黑格尔法哲学研究》，美茵河畔法兰克福，1969年，第42页以下。（见本书第85页以下。——译者）}

\section*{2}

如果要把握《法哲学原理》体系中自然和自由的关系以及《法哲学原理》在现代自然法终结处的独特的矛盾地位，那么我们就必须注意到黑格尔对霍布斯、卢梭和康德的法原则的这种评价。黑格尔在下述两种意义上站在霍布斯预备好的自然法理论的基地上：一方面，是在对于那种自然法理论具有构成性意义的自然概念的基地上；另一方面，则是在一种从既有的自然性和历史性权力中解放出来的意志的前提下，这种意志必须通过它自身的运动才与一切给定性发生关系，并且使这些给定的东西与自身和解。《法哲学原理》与17、18世纪自然法共享了一种自然的理念，这种自然本质上是否定性的，是通过消除“自然性”目的以及这些自然性目的对“最终”目的的隶属而得到界定的。作为精神哲学的一部分，它并不承认经院哲学中的那种自然的阶梯式进展。按照经院哲学的观点，这种阶梯式进展一直会进到“恩典的王国”(regnum gratiae)，其每一阶梯都有各自的自然地位。《法哲学原理》也不承认[莱布尼茨的]那种“自然王国”与“目的王国”之间的前定和谐学说：[按照这种学说，]“自然王国”与“目的王国”服从不同的原因和法则，并由于两者之间的这种和谐，自然的道路会从自身通向恩典。\footnote{莱布尼茨：《单子论》（1714年），第88节。}《法哲学原理》阐述的是“精神王国”的实现。“精神的王国”以一种目的论为前提，这种目的论内在于自由意志自身之中，并且就是渗透到一切实在性并在其中获得其自由的那种“概念”的运动。精神的阶梯式进展并不处于自然概念之下；它基于它自身的概念。而对黑格尔而言，这也就意味着：精神乃是基于自由之上。只有不断地摆脱自然，并且产生一种“第二自然”即精神的历史世界，自由才拥有其定在。精神的王国是自由的王国，它作为这样的王国不是自然性的等级制度，或者无世界的精神王国，而是“从它自身产生出来的精神的世界”。\footnote{《法哲学原理》，第4节；《哲学全书》，第381节。}

在这一世界——“第二自然”——中，自然与自由就像精神对自然那样互不相容。“精神的自然可以在其最完全的对立中认识到。我们把精神与物质对立起来。就像重量构成物质的实体一样，因此我们必须说，自由乃是精神的实体。”\footnote{《历史中的理性》，第5版，J. 荷夫迈斯特出版，汉堡，1955年，第55页；《哲学全书》，第381节，补充，《全集》第10卷，第19页以下。}在自然法理论中，同一种对立支配了对于自然法理论为构成性的自然的概念。这个概念在自然状态学说中扮演了一种既根本又不幸的角色。对黑格尔而言，这个概念具有一种“重要的、会导致绝对错误的歧义。一方面，自然意味着自然性存在，正如我们发现自己在不同的方面是直接造成的那样，是我们存在的直接方面。与之相对并与之相区分，自然亦是概念，事物的自然[本性]也就是事物的概念，就是事物的合乎理性方式的存在，这种事物完全不同于单纯自然性的东西”。\footnote{《历史中的理性》，第117页。}当“自然”指的是事物的概念时，那么自然状态中的法就必须被把握为那种“按照人的概念、按照精神的概念应当属于人的法”。然而，这不能与精神在其自然状态、在不自由与依赖于自然的状态中所是的东西混淆起来。因此，黑格尔与霍布斯、斯宾诺莎一起说道：必须走出自然状态(Exeundum est e statu naturae)。

在18世纪现代自然科学背景下，这句话的去目的论意义获得了一种独特的尖锐性。只要自然还是被规定为一种由目的所推动的秩序，那么自由就会与自然关联起来，并且人的行动也会被思维为对自然的延续或者模仿。随着自然被构建为一种可以以数学的方式加以描述和以力学的方式组织起来的因果关系，返回到这种自然就成为不可能的了；康德的《纯粹理性批判》为整个近代时期阐明了自由的思想，这种自由思想在自然中再也找不到类似于目的的东西，而是立足于自身之上。这对于黑格尔的《法哲学原理》具有决定性的意义。尽管黑格尔的《法哲学原理》也在重量概念中为其自由的基本概念找到一种自然式的类比，然而这里所涉及的是一个奠定了自从牛顿物理学以来系统理解自然的基础的范畴。《法哲学原理》第4节的补充指出，意志自由最好“参照”物理自然得到解释。自由和重量是意志和物体的“根本规定”。也就是说，这些规定对于意志和物体来说不是偶然的，而是必然的和本质的。它们的基本规定是：“当我们说物质是有重量的，我们可能认为这个谓语只是偶然的。然而并非如此，因为没有一种物质没有重量，其实物质就是重量本身。重量构成物体，而且就是物体。说到自由和意志也是一样，因为自由的东西就是意志。意志而没有自由，只是一句空话；同时，自由只有作为意志，作为主体，才是现实的。”\footnote{《法哲学原理》，第4节，补充，第50页。（中译文引自《法哲学原理》，第11-12页。——译者）}在重力中，自然物体追求一个中心点，也就是由自然物体在空间中的位置和关系所决定的物理体系的重心，但是这个重心对于这些自然物体而言始终是外在的中心。与之相反，在作为意志的基本规定的自由中，意志始终是自身相关的，这种自身相关正好就是意志的基础或者实体。意志尽管也追求一个中心点，但是追求的是一种它已经在自身中拥有的中心点：“它不在自身之外拥有统一；它始终在它自身中发现这种统一，它在自身中且依于自身；物质在它之外有其实体。与之相反，精神则是自在地存在，而这正是自由。”\footnote{《历史中的理性》，第55页。}

黑格尔与现代自然法所共享的第二个前提，涉及的是“自由意志”及其“概念”的起源。《法哲学原理》的基础正是“自由意志”及其“概念”的这种起源，而不是亚里士多德和经院哲学的自然法理论所刻画的那种具有复杂层级结构的统治和社会团体的目的论的阶梯式进展。尽管黑格尔明确地与霍布斯、卢梭和康德的理论保持距离，然而《法哲学原理》的整体安排在一个根本点上仍然是自然法式的：《法哲学原理》的概念发展始于与自然物和其他个体意志发生关系（财产和契约）的“一个主体的个体意志”\footnote{参见《哲学全书》，第400节以下；《法哲学》，第34节以下。——黑格尔一再强调这个属于自由意志的直接性的个别性环节，参见第13节附释、第39、43、46、52节。}，也就是始于“抽象”法的法，这种法重演了自然法学说的前政治状态。我们常常忽视了从个体意志出发的构造贯穿了整个体系，甚至延伸到对凝聚在“国家”中的意志的演绎。在黑格尔看来，这个意志\footnote{指凝聚在“国家”中的意志。——译者}也同样必须是一个“个体”的意志（君主的意志）。\footnote{参见《法哲学原理》，第279节，及附释：“任何科学的内在发展，即从它的简单概念到全部内容的推演[……]都显示出这样一个特点：同一个概念（在这里是意志）开始时（因为这是开始）是抽象的，它保存着自身，但是仅仅通过自身使自己的规定丰富起来，从而获得具体的内容。”（第381节）（中译文引自《法哲学原理》，第296页。——译者）另参见第190节，附释，以及《哲学全书》，第2版，第514节，《全集》第10卷，第397页。}

尽管如此，这个个体意志并不是自然法意义上的“自由”意志，并不是每个人的偶然的、个别的任意；相反，法、普遍意志的概念正是产生于对每个人的偶然的、个别的任意的交互限制。因为只有意志概念在主观精神学说中从一切偶然的自然规定性中被纯化之后，黑格尔才引入作为《法哲学原理》的对象的意志概念，所以《法哲学原理》的概念起源就是从包含在个体意志中的“普遍意志”出发的。因此，对黑格尔而言，自由意志是“理性意志和个体意志的统一，个体意志是理性意志的活动的直接的和独特的要素”。\footnote{参见《哲学全书》，第2版，第485节，第382页以下；另参见《法哲学原理》，第2节，第39页；第4节，第50页以下。}然而，这并不意味着取消自然法思想，实际上毋宁是对它的一个巨大的加强。\footnote{参见罗森茨威格：《黑格尔与国家》第2卷，慕尼黑/柏林，1920年，第106页以下。}借此，个体不再作为其自身，在其自然性中，而是作为理性存在者，而成为法学的出发点和对象。当康德和费希特把他们所发现的自由概念局限在伦理学领域并把自然法建立在每个人的（由自然所规定的）自由任意之上时，黑格尔也把这种自由概念运用于法学中。法不限制自由意志，而是自由意志的“定在”，即“作为理念的自由”。对黑格尔来说，这一措辞说的是：自由不再仅仅是规范和公设（康德意义上的理念），而是事实，它存在于社会历史世界之中，而不仅仅是某种没什么指望的东西。

为了思维这种理性的事实，就需要黑格尔在《法哲学原理》导言中将自然法的意志“概念”应用于其下的那种辩证法。依照黑格尔的观点，康德尽管在先验哲学的自我原则中发现了意志的自我决定的奠基性环节，即“纯粹的无规定性或自我在自身中的纯粹反思”的环节，在其中，所有的限制和每一种由自然“给予的”内容都消解无余了，但是自为地看，意志的自由仿佛只是消极的；在“政治”中，意志自由（如同在以卢梭理论为范本的法国大革命进程中所看到的那样）导致了“一切现存社会秩序的破坏”和“对任何企图重整旗鼓的组织的消灭”。\footnote{《法哲学原理》，第5节，第54页以下。（中译文参见《法哲学原理》，第14页。——译者）}因此，自我作为意志的纯粹概念，必然在自身中包含了第二个环节，即“规定和设定一种规定性作为一种内容和对象”，借此自我将自身设定为有规定的东西，也就是说进入定在。这种对意志予以规定和限制的特殊物，绝不是任何给定的外在限制，而是（在黑格尔看来，康德和费希特忽略了这一点）内在于自我决定的行为之中的。在这两个环节（即意志的无规定的普遍性和有规定的特殊性）的统一中，作为意志的“概念”的自由将自身实现为第三个环节，——意志在他在中自在地存在，或者说“自我的自身规定，在一中，它设定自己为[……]被限制的东西；它留在自己那里，即留在与自己的同一性和普遍性中；又它在这一规定中只与自己本身联结在一起”。\footnote{同上书，第7节，第58页以下。（中译文引自《法哲学原理》，第17页。译文稍有改动。——译者）}

\section*{3}

自由意志的辩证运动将普遍性的环节与特殊物\footnote{指普遍性的环节与特殊物。——译者}中介起来，也将这两者与它自身中介起来，而成为“具体的”个体性、“自由意志的定在”（《法哲学原理》，第29节）。自由意志的这一辩证运动，就是我们在黑格尔《法哲学原理》中作为“概念”而见到的那种“事物的自然[本性]”。概念把自然的意义完全把握在自身中；概念指的无非就是把它自身以及它的所有规定都建立在“自由的人格性”上面的法，也就是“一种毋宁说是自然规定的对立面的自我规定”。\footnote{《法哲学原理》，第7节，第58页以下。（里德尔原注疑误。该书第7节没有这段引文。——译者）也参见《纽伦堡哲学入门》，《全集》第3卷，第55、71页。特别是《纽伦堡哲学入门》，第三篇，第2章；《哲学全书》，第三部，二之1 法，第181节：“精神作为自由的、具有自我意识的存在者，是自相等同的自我，这种自我在其绝对否定关系中首先是排斥性的自我，即个别的、自由的存在者，或人格(Person)。”[本段引文见《黑格尔全集》（历史考订版）第10卷第1分册，第355页。中译文引自《黑格尔全集》第10卷，第289页。译文稍有改动。——译者]}概念的辩证运动开始于摆脱了所有“自然性”规定的“人格”的自我规定，也就是开始于作为法的唯一原则的自由意志的无限的、不受任何外物限制的“概念”。黑格尔在柏林时期的法哲学讲座中反复讲述了概念（即自然）、自由和法之间的这种联系。第一次柏林法哲学讲座在1818-1819年、《法哲学原理》出版之前。这次讲座第3节说道：“法的原则不在于自然，既不在外在自然，也不在人的主观自然，因此时人的意志是自然地被规定的，也就是说，是欲望、冲动和偏好的领域。法的领域是自由的领域，在其中，自由外在化自身并赋予自身以实存，此时尽管自然会出场，但却是作为一种没有独立性的东西[而出场的]。”\footnote{《黑格尔教授1818-1819年冬季学期讲授的自然法和国家法》（Natur- und Staatsrecht nach d. Vortrag des Professors Hegel im Winterhalbjahr 1818/1819），霍迈尔(G. Homeyer)笔记，四开本德语手稿第1155号，第9页。[本段引文见《黑格尔全集》（历史考订版）第26卷第1分册，第239页。——译者]}正如法哲学所述，从抽象法的现象形式开始，直到家庭、市民社会和国家为止，自由的各种外化就是社会制度和行为方式的所有环节。毫无疑问，在所有这些环节中，处处都有“自然性”联系的身影。但问题在于，这些“自然性”联系是否从自身产生出一种支配那些形式的法律。在黑格尔看来，必然性——比如说，在与个体生命的关系中，国家的定在所具有的那种必然性——不再意味着，对于个体而言，必须在国家中生活乃是一种自然的规律。毋宁说，国家的必然性建立在自由自身所立的法律基础上。黑格尔对第3节的解说毫不含糊地指出：“这一节是由自然法之名引发的。自由与必然性的统一并不是通过自然，而是通过自由产生出来的。自然事物一成不变，自己并不能从规律中摆脱出来，以便自己制定法律。然而，精神则使自身从自然中挣脱出来，并且自己产生它的自然、它的法则本身。因此，自然不是法的基地。”与康德和费希特的看法一样\footnote{参见康德：《反思录》，第7084条，康德：《全集》（科学院版）第19卷，第245页。费希特：《自然法权基础》，《全集》第3卷，第148页。}，黑格尔也认为自然法的名称“不过是习传下来的”，按照事物[本身]、法的“概念”，这一名称甚至是不正确的，因为“自然可以理解为：(1)本质、概念。(2)无意识的自然（真正的意义）”；“真正的名称”必须是“哲学法学”。\footnote{《自然法和国家法》，霍迈尔笔记，第9页。[本段引文见《黑格尔全集》（历史考订版）第26卷第1分册，第240页。——译者]}

因此，在起草法哲学体系的时期，黑格尔的思想发展完全回到了《论自然法的科学探讨方式》一文想要通过“伦理自然”法的建构予以克服的那种自由哲学的基础上。在结束了对“哲学法学”基本概念的漫长理解过程（这一过程同时也澄清了哲学法学与传统自然法的历史关系）之后，现在对此也就无须赘言了。在黑格尔的体系中，自然与自由、自然规律与法权法则分离开来了。因而黑格尔决然地拒绝传统观念：就只有人可以认识和遵守自然规律而言，“自然规律”对人来说可以成为法的“范本”；人的历史性定在所唯一具有的法则，就是自由法则，不是“自然”，而是“概念”自身制定了这一法则。\footnote{参见上书，第3页：“对我们而言，事物的概念并不来自自然。”[本段引文见《黑格尔全集》（历史考订版）第26卷第2分册，第774页。——译者]}概念的自我立法——对于黑格尔而言，这意味着18世纪发现的人的普遍法权能力的观念\footnote{参见上书，第5页：“这须被尊重为某种伟大之事：现在人因其为人，必须认为拥有权利，以至于他的人的存在要高于他的身份。”}，哲学通过这种观念去反对现成的东西，并且打破了自然性的假象，在这种假象背后，“法的思想”、自由都还一直隐而不彰。这种思想将现成的法和对它的知识变为同一个东西。黑格尔在1824-1825年的讲座导言中说，尽管在以自然、自然需要和目的等为基础的“习惯的自然法”中，自由的思想并没有被有意排除掉，但“事实上，由于没有领会到这两种原则的独特性就接受了它们，从而怠慢了自由”。自由不能作为自然，而是必须在“概念”的规范性中去思维，在自由与其自身以及自由与自然之间，概念起着中介的作用。当自由以“概念”的要求出现时，这种要求就包含了“自由仅仅指与法和伦理有关的自由，这样自然法之名就陷入了动摇之中。”\footnote{《法哲学讲演录》，格里斯海姆笔记，第9页以下。[本段引文见《黑格尔全集》（历史考订版）第26卷第3分册，第1054页。——译者]}只有当我们“承认”，自由就是自然法所公设的“事物的自然[本性]”——也就是“概念”——的时候，我们才能保住自由的概念。然而，这种表达是“不恰当的”，因为不同于自由，自然“是无拘束的，并不与某个它物相对立。而自由则相反，同时显得是好战的，并且具有众多的对立物，而它的头一个对立物就是自然本身”。\footnote{同上书，第10页。[本段引文见《黑格尔全集》（历史考订版）第26卷第3分册，第1054页。——译者]}

\section*{4}

对于黑格尔法哲学的体系安排和概念处理来说，自然与自由的对立具有决定性的意义。这一对立不仅决定了这部著作两个标题中“哲学法学”和“自然法”的区分\footnote{黑格尔在1824-1825年讲座中明确提到了这一点。参见格里斯海姆笔记，第3页；以及拙文“黑格尔《法哲学原理》中的传统与革命”见《黑格尔法哲学研究》，前揭，第100页以下。（见本书第170页以下。——译者）}，而且也决定了体系最后也是最重要的部分随着在家庭与国家之侧引入“市民社会”概念而经验到的双重化，因为区分市民社会和国家的体系根据在于自然概念和自由概念的分离。黑格尔在历史中思考终结这种分离。在近代自然法的发展进程中，国家(civitas, res publica)与市民社会(societas civilis)的同一瓦解了。这种同一曾在传统的“封建”国家和社会理论上面打下烙印，并且在哲学上通过自然规律(lex naturalis)的目的论而得到合法化：这种目的论让个体的“自然性目的”（冲动、需要等）与整体的伦理目的彼此协调一致。然而，这种同一性的瓦解并没有导致一种持久的差别，而是再次被克服了。吊诡的是，其原因恰恰在于自然法的契约理论。这种理论既是分离的催化剂，同时又是克服这种分离的手段。一方面，自然的目的论的、超出个体和市民社会之定在的运动秩序遭到了抛弃，而另一方面，自然又在自然法学说的出发点上，作为多的“任意”的环节，进入契约的建构之中。如果每个人的任意都应当能够按照一条普遍的法则（这一普遍法则取代了被废弃的自然法则）而与所有人的任意相一致的话，那么法概念所必须限制的，就正是包含在“自由任意”这一说法之中的那些需要、冲动和爱好。或许在这里可以找到黑格尔之所以反对康德和卢梭之引起争议的法原则与自由原则的界分的真正动因，对此，黑格尔在《法哲学原理》中也非常重视。与此同时，黑格尔自己也对卢梭—康德与旧自然法之间在自然观上所具有的区别非常清楚。传统的探讨方式（在这种探讨方式中，自然概念的“两种含义”还没有发生分离）以人的“自然性存在”即“爱好、需要、冲动”为基础：首先是肉体的需要及其满足的必然性，然后是“社交的冲动、社会的冲动，这种冲动一方面涉及两性关系，另一方面也扩展到普遍的市民社会”，最后自身“发展为国家”。\footnote{参见《1824-1825年讲演录》，格里斯海姆笔记，第5-8页。《法哲学原理》导论，第19节，附释。}而对黑格尔而言，这里正好表现了“这种事情的不足”。并不存在自然的目的论，可以对个人的冲动和需要的本性与“国家”的伦理历史性定在进行中介；相对于国家和市民社会按照“概念”之所是，所有从中推演出来的规定都是“抽象的”：

\begin{quote}
如果我们必须说，市民社会、国家是从一种冲动中产生出来的，这马上就表明了从一种冲动到国家这种事情的不足。之前在哲学中被假定为国家之基础的社交冲动，乃是某种缺乏规定的、抽象的东西，对于具有丰富层次的国家来说，它只能提供少量的规定，而且对它所要奠基的东西来说，它也显得太过贫乏了。\footnote{参见格里斯海姆笔记，第8页。[本段引文见《黑格尔全集》（历史考订版）第26卷第3分册，第1053页。——译者]}
\end{quote}

霍布斯最先抛弃了这种自然法原则，其后卢梭与康德费希特哲学又把它与自由原则对立起来。但是，正如黑格尔在1824-1825讲座中所言，[卢梭、康德和费希特]“所采取的方式并非我们的方式，也没有让法学成为一种全面的和一以贯之的科学”。\footnote{同上书，第11页。〔本段引文见《黑格尔全集》（历史考订版）第26卷第3分册，第1054页。其中“我们的”在历史考订版中为“真正的”。——译者〕也可参见《1823-1824年讲演录》（Vorlesungen von 1823/1824），霍托笔记，第7页以下。}黑格尔同意康德，认为在自然规律之下自由是不可能的，因为自然和自由彼此对立。因此他和康德一样，拒绝传统目的论自然观中所包含的对“国家的历史起源”\footnote{参见《法哲学原理》，第258节，第330页。（中译文参见《法哲学原理》，第254页。——译者）}的追问，或者用政治学和自然法的传统术语来说，对“市民社会的开端”的追问，因为法哲学只是涉及“国家的理念”、“被思维的概念”。而在康德和费希特那里（以及他们之前的卢梭那里），他们作为法权学说基础的体系原则是与他们对体系原则所做的历史阐述相矛盾的。因为“思维本身”是作为意志，国家的原则“不仅在形式上（诸如社会性冲动、神的权威），而且内容上也是思想”，因此一方面这种原则把自然从自身中排除了出去，而另一方面又通过他们把仅仅在历史规定的形式上的“个体意志”中的作为“特殊个体”的意志置于顶端，而将自然重新引了进来。\footnote{参见《法哲学原理》，第258节，以及第29节，第79页以下，第182节，补充；《历史中的理性》，第111页以下。}结果，他们就不是把“普遍意志”、国家的“概念”理解为“意志的自在自为的合乎理性的东西，而是理解为从这种个体意志中产生出来的共同的东西”。由此，在国家和市民社会的“概念”方面，近代自然法和启蒙的古典自由主义又一起重新落到了之前的自然法学说的立场上。霍布斯、卢梭和康德作为前提来瓦解国家(civitas)和市民社会(societas civilis)的古典同一性的契约建构重复了那种对两个概念领域的语言体系上的“混淆”，在《法哲学原理》第182节中，黑格尔拿这种混淆去反驳“国家”的自然法“看法”。\footnote{同上书，第182节，补充：“如果把国家想象为各个不同的人的统一，亦即仅仅是共同性的统一，其所想象的只是指市民社会的规定而言。许多现代的国家法学者都不能对国家提出除此之外任何其他看法。”（第262页以下）（中译文引自《法哲学原理》，第254页。——译者）}

这一反驳之所以可能，前提是黑格尔自己从近代自然法的出发点所得出的结论。通过把自然规律和自由法则彻底对立起来，黑格尔不仅抛弃了传统自然法名称，而且也放弃了市民社会(societas civilis)这一说法。与抛弃传统自然法概念相比，后者所造成的影响要深远得多。随着黑格尔引入国家和市民社会的分离，一个问题视域被打开了。在这一问题视域之内，各个个体的生命和意志的关系、他们彼此之间的联系以及他们与自然的联系就能在理论上获得独立且相对隔绝开来，作为“社会”联系得到研究和把握。恰恰因为个体意志的自然规定性、“特殊的个体”及其需要、冲动和爱好与普遍的、在国家中自身显现着的意志的演绎并没有任何直接的论证关系，所以黑格尔才有可能发现那种作为现代市民社会的经济基础的“需要的体系”。而对自然法理论的自由主义信徒们来说，这个需要的体系却被社会契约模式及其政治意涵遮蔽了。当然，在黑格尔这里，自然和自由的关系问题在《法哲学原理》中也并没有得到满意的解决，而是在一个新的层次上重新出现了。“作为理念的自由”在概念上要求市民社会从国家中脱离出来，由此，自由也就在[市民社会]这一层面从现实性坠入表面现象、“伦理的现象世界”。在这个世界中，伦理的同一性不能作为自由，只能作为必然性而得到思考。\footnote{《法哲学原理》，第186节，第267页；参见第184-185节。}这里，集中于自由意志领域的概念运动已经蓄势待发，准备转变为一种社会的运动法则，这种运动法则的基础不是作为概念之前提的自由，而是自然——“不平等的要素”。“不平等的要素”表达了精神的特殊性的客观法。其中，不仅自然设定的不平等没有被扬弃，反而“从精神中产生出来”。\footnote{同上书，第200节，附释，第278页以下。（中译文参见《法哲学原理》，第211页。——译者）此外，黑格尔在对世袭君主进行演绎（第280节以下）和对战争进行辩护（第324节）时，都援引了概念所无法化解的自然作为依据。}这里构成了突破黑格尔法哲学的地方之一。在这些突破黑格尔法哲学的地方，概念的辩证起源达到了它的界限——也就是黑格尔的最重要的、最具有历史影响力的学生马克思想要通过把这种辩证法“唯物主义地颠倒”为一种生产力和生产关系的运动而加以超越的那种界限。